% 泊松代数（综述）
% license CCBYSA3
% type Wiki

本文根据 CC-BY-SA 协议转载翻译自维基百科\href{https://en.wikipedia.org/wiki/Poisson_algebra}{相关文章}

在数学中，泊松代数是一种带有李括号的结合代数，并且该李括号满足莱布尼茨法则；也就是说，括号运算同时是一个导子。泊松代数在哈密顿力学中自然出现，同时在量子群研究中也处于核心地位。带有泊松代数结构的流形称为泊松流形，其中辛流形和泊松–李群是特殊情形。该代数的名称是为了纪念西美昂·德尼·泊松。
\subsection{定义}
一个泊松代数是定义在域$K$上的向量空间，配备了两个双线性运算$\cdot$和$\{,\}$，并满足以下性质：
\begin{itemize}
\item 运算 $\cdot$ 形成一个结合的 $K$-代数；
\item 运算 $\{,\}$，称为泊松括号，形成一个李代数，因此它是反对称的，并且满足Jacobi 恒等式；
\item 泊松括号对结合积 $\cdot$ 起导子作用，即对代数中任意三个元素 $x, y, z$，有$\{x, y \cdot z\} = \{x, y\} \cdot z + y \cdot \{x, z\}$.
\end{itemize}
最后一个性质通常允许对该代数给出多种不同的表述，下面的例子中会有所体现。
\subsection{例子}
泊松代数出现在多种情境中。
\subsubsection{辛流形}
在一个辛流形上，所有取实值的光滑函数空间构成一个泊松代数。在辛流形上，每一个实值函数 $H$ 都会诱导一个向量场$X_H$，称为哈密顿向量场。于是，对于任意两个定义在辛流形上的光滑函数$F$和$G$，泊松括号可定义为：
$$
\{F,G\} = dG(X_F) = X_F(G).~
$$
这个定义部分上之所以一致，是因为泊松括号起到导子的作用。等价地，可以将括号 $\{,\}$ 定义为：
$$
X_{\{F,G\}} = [X_F, X_G],~
$$
其中 $[,]$ 表示李导数。

当辛流形是带有标准辛结构的 $\mathbb{R}^{2n}$ 时，泊松括号取著名的形式：
$$
\{F,G\} = \sum_{i=1}^n 
\frac{\partial F}{\partial q_i}\frac{\partial G}{\partial p_i} - 
\frac{\partial F}{\partial p_i}\frac{\partial G}{\partial q_i}.~
$$
类似的思想也适用于泊松流形，它通过允许辛双向量退化来推广辛流形的概念。
\subsubsection{李代数}
李代数的张量代数具有泊松代数结构。在关于普遍包络代数的文章中可以找到一个非常具体的构造。

构造过程如下：首先构造该李代数底层向量空间的张量代数。张量代数就是该向量空间所有张量积的“不交并”（即直和$\oplus$）。然后可以证明，李括号能够一致地提升到整个张量代数上：它满足乘积法则以及泊松括号的 Jacobi 恒等式，因此当被提升时，它就是泊松括号。此时，运算对$\{,\}$与$\otimes$便形成了一个泊松代数。需要注意的是，$\otimes$既不是交换的，也不是反交换的；它只是结合的。

因此，可以得出一个一般性的结论：任意李代数的张量代数都是泊松代数。而普遍包络代数则是通过对该泊松代数结构取商而得到的。
\subsubsection{结合代数}
若$A$是一个结合代数，那么通过施加对易子$[x,y] = xy - yx$可以将其转化为一个泊松代数（因此也是一个李代数），记作 $A^L$。需要注意的是，这里得到的 $A^L$ 不应与上一节所述的张量代数构造混淆。如果愿意，也可以对其应用张量代数的构造，但那会得到一个不同的泊松代数，而且规模更大。
\subsubsection{顶点算子代数}
对于一个顶点算子代数$(V, Y, \omega, 1)$，其商空间$V / C_2(V)$是一个泊松代数，其中$\{a, b\} = a_0 b$，并且 $a \cdot b = a_{-1} b$。对于某些顶点算子代数，这些泊松代数是有限维的。
\subsubsection{ \(Z_2\) 分次}
泊松代数可以通过两种不同的方式赋予$\mathbb{Z}_2$-分次。这两种方式分别得到泊松超代数和Gerstenhaber 代数。它们的区别在于括号对分次的影响。

对于泊松超代数，分次由下式给出：
$$
|\{a,b\}| = |a| + |b|,~
$$
而在 Gerstenhaber 代数中，括号会使分次降低 1：
$$
|\{a,b\}| = |a| + |b| - 1.~
$$
在这两种表达式中，$|a| = \deg a$表示元素$a$的分次；通常，这个分次刻画了 $a$ 是如何由生成元的偶积或奇积分解而成的。Gerstenhaber 代数通常出现在 BRST 量子化中。
